\documentclass[12pt,a4paper]{article}
\usepackage[T1]{fontenc}
\usepackage[utf8]{inputenc}
\usepackage[provide=*,lithuanian]{babel}
\usepackage{amsmath,amssymb}
\usepackage[margin=2cm]{geometry}
\usepackage{enumitem,multicol}
\setlength{\parindent}{0pt}
\pagestyle{empty}
\setlist[enumerate,1]{label=\arabic*.,leftmargin=*,itemsep=0.65em}
\setlist[enumerate,2]{label=\alph*),leftmargin=*,itemsep=0.6em}

\begin{document}
\begin{center}
  {\Large Namų darbas}\\[4pt]
  {\large Kvadratinis trinaris}
\end{center}

\begin{enumerate}
\item Išskaidykite dauginamaisiais.
\begin{multicols}{2}
\begin{enumerate}
  \item $x^2-7x+12$;
  \item $2x^2+7x+3$;
  \item $3x^2-5x-2$;
  \item $-2x^2+5x+3$.
\end{enumerate}
\end{multicols}

\item Nustatykite kiekvieno kvadratinio trinario ekstremumo tašką ir ekstremumą. Nurodykite, ar rastasis ekstremumas yra minimumas, ar maksimumas.
\begin{multicols}{2}
\begin{enumerate}
  \item $P(x)=x^2-6x+5$;
  \item $P(x)=-2x^2+8x-3$;
  \item $P(x)=3x^2+12x+7$;
  \item $P(x)=-x^2-6x+2$.
\end{enumerate}
\end{multicols}

\item Įrodykite šiuos teiginius. Kiekvienu atveju nurodykite, su kuria $x$ reikšme pasiekiama nurodytoji reikšmė.
\begin{enumerate}
  \item Daugianario $P(x)=-3(x+2)^2+8$ didžiausia reikšmė yra $8$.
  \item Daugianario $P(x)=6(x-1)^2-7$ mažiausia reikšmė yra $-7$.
\end{enumerate}

\item Išskirkite pilnąjį kvadratą.
\begin{multicols}{2}
\begin{enumerate}
  \item $x^2+8x+11$;
  \item $2x^2-12x+7$;
  \item $-3x^2-12x+5$;
  \item $4x^2+8x-1$.
\end{enumerate}
\end{multicols}

\item Užrašykite kvadratinį daugianarį $P(x)$, kurio ekstremumo taškas lygus $2$, ekstremumas lygus $4$, o $P(1)=1$.

\item Užrašykite antrojo laipsnio daugianarį $P(x)$, kurio vyriausiasis koeficientas lygus $2$, $P(1)=P(7)$, o ekstremumas lygus $-2$.

\item Duotas antrojo laipsnio daugianaris $P(x)$. Žinoma, kad $P(4)=P(16)$ ir $P(7)=1$. Apskaičiuokite $P(13)$.

\item Kvadratinio trinario $P(x)$ vyriausiasis koeficientas yra neigiamas ir $P(-\pi)=P(\pi)$. Išrikiuokite skaičius $P(1)$, $P(-10)$, $P(6)$ ir $P(-5)$ mažėjimo tvarka. Atsakymą pagrįskite.

\item Kvadratinio daugianario $P(x)$ diskriminantas lygus $16$, vyriausiasis koeficientas yra teigiamas, o ekstremumo taškas lygus $1$. Raskite daugianarį $P(x) $, jei viena iš jo šaknų yra $3$.

\item Kvadratinio trinario diskriminantas lygus $100$, ekstremumas lygus $5$, o $P(0)=5$. Raskite trinarį $P(x)$  ir nustatykite, ar ekstremumas yra minimumas, ar maksimumas.

\item Kvadratinio daugianario $P(x)=2x^2+bx+c$ ekstremumo taškas lygus $-1$, o $P(-1)=-2$. Raskite šio daugianario šaknų kvadratų sumą.
\end{enumerate}

\section*{Atsakymai}
\begin{enumerate}
\item
\begin{enumerate}
  \item $(x-3)(x-4)$;
  \item $(2x+1)(x+3)$;
  \item $(3x+1)(x-2)$;
  \item $(-2x-1)(x-3)$.
\end{enumerate}

\item
\begin{enumerate}
  \item Ekstremumo taškas $x=3$, minimumas $-4$;
  \item Ekstremumo taškas $x=2$, maksimumas $5$;
  \item Ekstremumo taškas $x=-2$, minimumas $-5$;
  \item Ekstremumo taškas $x=-3$, maksimumas $11$.
\end{enumerate}

\item
\begin{enumerate}
  \item Didžiausia reikšmė $8$, kai $x=-2$;
  \item Mažiausia reikšmė $-7$, kai $x=1$.
\end{enumerate}

\item
\begin{enumerate}
  \item $(x+4)^2-5$;
  \item $2(x-3)^2-11$;
  \item $-3(x+2)^2+17$;
  \item $4(x+1)^2-5$.
\end{enumerate}

\item $P(x)=-3(x-2)^2+4$.

\item $P(x)=2(x-4)^2-2$.

\item $P(13)=1$.

\item $P(1)>P(-5)>P(6)>P(-10)$.

\item $P(x)=(x-1)^2-4$.

\item $P(x)=-5x^2+5$; maksimumas $5$.

\item $4$.
\end{enumerate}
\end{document}
